% =====================================================================
% Trabajo1.tex (standalone) — COMPILA CON pdfLaTeX (NO fontspec / NO minted)
% =====================================================================
\documentclass[12pt]{article}

% ---------- Idioma / codificación (pdfLaTeX) ----------
\usepackage[utf8]{inputenc}
\usepackage[T1]{fontenc}
\usepackage[spanish]{babel}

% ---------- Formato ----------
\usepackage[a4paper,margin=2.5cm]{geometry}
\usepackage{microtype}
\usepackage{parskip} % evita indentación y mejora espaciado

% ---------- Links ----------
\usepackage{hyperref}
\usepackage{xurl}
\hypersetup{
  colorlinks=true,
  linkcolor=black,
  urlcolor=blue,
  citecolor=black
}

% ---------- Figuras ----------
\usepackage{graphicx}
\usepackage{float}
\usepackage{caption}

% ---------- Código (sin shell-escape) ----------
\usepackage{listings}
\usepackage{xcolor}

\lstdefinestyle{mystyle}{
  basicstyle=\ttfamily\small,
  columns=fullflexible,
  breaklines=true,
  breakatwhitespace=false,
  frame=single,
  rulecolor=\color{black!30},
  tabsize=2,
  showstringspaces=false,
  upquote=true
}
\lstdefinestyle{cstyle}{
  style=mystyle,
  language=C
}
\lstdefinestyle{txtstyle}{
  style=mystyle
}

% ---------- Utilidades ----------
\newcommand{\HRule}{\par\noindent\rule{0.5\linewidth}{0.5pt}\par}

% Control de tamaño para que NUNCA se salga de la página:
\newcommand{\fitimage}[2][]{%
  \begin{center}
    \includegraphics[
      width=\linewidth,
      height=0.83\textheight,
      keepaspectratio,
      #1
    ]{#2}
  \end{center}
}

\begin{document}

\section{Resolución del Problema de los Filósofos Comensales usando Semáforos}
\label{sec:filosofos-semaforos}

\textbf{Curso:} Sistemas Operativos

\textbf{Docente:} Fabio Andres Guzman Figueroa

\textbf{Trabajo 1:} Sincronización con Semáforos

\textbf{Estudiante:} Santiago Hernández Morantes -- Simón Bedoya Urrea

\textbf{Fecha:} 25 de febrero de 2026

\HRule

\newpage

\subsection{Introducción}
\label{subsec:introduccion}

En sistemas concurrentes, múltiples procesos o hilos pueden competir por recursos compartidos.
Cuando no existe un mecanismo de sincronización adecuado, pueden presentarse problemas como:

\begin{itemize}
  \item Condiciones de carrera
  \item Interbloqueo (\textit{deadlock})
  \item Inanición (\textit{starvation})
\end{itemize}

Los semáforos, propuestos por Edsger Dijkstra, son una herramienta fundamental para controlar el acceso a recursos compartidos y evitar estos problemas.

En este trabajo se implementa el clásico problema de los Filósofos Comensales utilizando semáforos POSIX en C.

\subsection{Descripción del Problema}
\label{subsec:descripcion}

El problema plantea:

\begin{itemize}
  \item $N$ filósofos sentados alrededor de una mesa circular.
  \item Entre cada par de filósofos hay un cubierto.
  \item Para comer, un filósofo necesita dos cubiertos: el de su izquierda y el de su derecha.
\end{itemize}

Cada filósofo alterna entre:

\begin{itemize}
  \item Pensar
  \item Tener hambre
  \item Comer
\end{itemize}

El reto es garantizar que:

\begin{itemize}
  \item No haya interbloqueo (\textit{deadlock}).
  \item No haya uso simultáneo de un mismo cubierto.
  \item El sistema sea concurrente (no estrictamente secuencial).
\end{itemize}

\subsection{Condiciones del Deadlock}
\label{subsec:condiciones-deadlock}

Para que ocurra un deadlock deben cumplirse cuatro condiciones:

\begin{enumerate}
  \item \textbf{Exclusión mutua:} los cubiertos no pueden compartirse.
  \item \textbf{Hold and wait:} un filósofo puede tener un cubierto y esperar otro.
  \item \textbf{No apropiación:} no se puede forzar a soltar un cubierto.
  \item \textbf{Espera circular:} cada filósofo espera el recurso del siguiente.
\end{enumerate}

Nuestra solución rompe la condición de \textbf{espera circular} mediante una verificación controlada del estado de los vecinos antes de habilitar el estado de comida.

\subsection{Estrategia de Solución}
\label{subsec:estrategia}

Se implementó la solución clásica de Dijkstra basada en:

\begin{itemize}
  \item Un semáforo global \texttt{mutex} para proteger la sección crítica.
  \item Un semáforo privado \texttt{s[i]} por cada filósofo.
  \item Un arreglo \texttt{state[i]} con estados: \texttt{THINKING}, \texttt{HUNGRY}, \texttt{EATING}.
\end{itemize}

\textbf{Idea principal:} un filósofo solo puede comer si está hambriento y sus dos vecinos no están comiendo. Si no cumple estas condiciones, queda bloqueado en su semáforo privado hasta que pueda continuar.

Esto permite:
\begin{itemize}
  \item Evitar interbloqueo
  \item Maximizar concurrencia (no vecinos pueden comer en paralelo)
  \item Asegurar exclusión mutua sobre los cubiertos
\end{itemize}

\subsection{Esquema del Algoritmo}
\label{subsec:esquema}

\textbf{Estados:} \texttt{THINKING} $\rightarrow$ \texttt{HUNGRY} $\rightarrow$ \texttt{EATING} $\rightarrow$ \texttt{THINKING}

\textbf{Funciones clave:}
\begin{itemize}
  \item \texttt{take\_forks(i)}
  \item \texttt{put\_forks(i)}
  \item \texttt{test(i)}
\end{itemize}

\subsection{Pseudocódigo Explicado}
\label{subsec:pseudocodigo}

\lstset{style=txtstyle}
\begin{lstlisting}
PSEUDOCODIGO — FILOSOFOS COMENSALES (Dijkstra con semaforos)

CONSTANTES:
    N
    THINKING <- 0
    HUNGRY   <- 1
    EATING   <- 2

FUNCIONES AUXILIARES:
    LEFT(i)  <- (i + N - 1) mod N
    RIGHT(i) <- (i + 1) mod N

VARIABLES COMPARTIDAS:
    state[0..N-1]
    semaphore mutex
    semaphore s[0..N-1]

PROCEDIMIENTO test(i):
    if state[i] == HUNGRY and
       state[LEFT(i)]  != EATING and
       state[RIGHT(i)] != EATING then

        state[i] <- EATING
        signal(s[i])
    end if

PROCEDIMIENTO take_forks(i):
    wait(mutex)
        state[i] <- HUNGRY
        test(i)
    signal(mutex)

    wait(s[i])   // se bloquea si no puede comer
    // aqui ya esta en EATING

PROCEDIMIENTO put_forks(i):
    wait(mutex)
        state[i] <- THINKING
        test(LEFT(i))
        test(RIGHT(i))
    signal(mutex)

RUTINA philosopher(i):
    repetir ROUNDS veces:
        pensar
        take_forks(i)
        comer
        put_forks(i)
\end{lstlisting}

\subsection{Justificación de Corrección}
\label{subsec:correctitud}

\textbf{Exclusión mutua de cubiertos:} un filósofo solo entra en estado \texttt{EATING} si ambos vecinos \textbf{no} están comiendo. Por tanto, no hay uso simultáneo de un mismo cubierto.

\textbf{No hay deadlock:} la habilitación a comer se hace mediante \texttt{test(i)} bajo protección de \texttt{mutex}, evitando la espera circular: o queda habilitado completo, o queda bloqueado.

\textbf{Concurrencia real:} filósofos no vecinos pueden comer simultáneamente, pues sus condiciones no interfieren.

\subsection{Implementación en C}
\label{subsec:implementacion-c}

Repositorio:
\href{https://github.com/Titolin4612/PR-SistemasOperativos/blob/main/Trabajos/trabajo1/trabajo1.c}{\texttt{GitHub: trabajo1.c}}

\lstset{style=cstyle}
\begin{lstlisting}
#define _DEFAULT_SOURCE    // expone usleep/useconds_t con -std=c11 -pedantic
#include <pthread.h>
#include <semaphore.h>
#include <stdio.h>
#include <stdlib.h>
#include <unistd.h>
#include <time.h>

#ifndef N_PHILOS
#define N_PHILOS 5
#endif

#ifndef ROUNDS
#define ROUNDS 1
#endif

#ifndef THINK_US
#define THINK_US 700000
#endif

#ifndef EAT_US
#define EAT_US 700000
#endif

typedef enum { THINKING, HUNGRY, EATING } state_t;

static state_t state[N_PHILOS];
static sem_t mutex;
static sem_t s[N_PHILOS];

static pthread_t th[N_PHILOS];
static int ids[N_PHILOS];

static struct timespec t0;

static inline int LEFT(int i)  { return (i + N_PHILOS - 1) % N_PHILOS; }
static inline int RIGHT(int i) { return (i + 1) % N_PHILOS; }

static long now_ms_since_start(void) {
    struct timespec t;
    clock_gettime(CLOCK_MONOTONIC, &t);
    long sec  = (long)(t.tv_sec  - t0.tv_sec);
    long nsec = (long)(t.tv_nsec - t0.tv_nsec);
    return sec * 1000L + nsec / 1000000L;
}

static void log_event(int i, int round, const char *msg) {
    printf("[%6ld ms] F[%d] R[%d/%d] %s\n",
           now_ms_since_start(), i + 1, round, ROUNDS, msg);
    fflush(stdout);
}

static void test(int i) {
    if (state[i] == HUNGRY &&
        state[LEFT(i)]  != EATING &&
        state[RIGHT(i)] != EATING) {

        state[i] = EATING;
        sem_post(&s[i]);
    }
}

static void take_forks(int i, int round) {
    sem_wait(&mutex);

    state[i] = HUNGRY;
    log_event(i, round, "hambriento...");
    test(i);

    sem_post(&mutex);

    sem_wait(&s[i]);

    char buf[128];
    snprintf(buf, sizeof(buf), "usa cubiertos (%d, %d) y come...",
             i + 1, RIGHT(i) + 1);
    log_event(i, round, buf);
}

static void put_forks(int i, int round) {
    sem_wait(&mutex);

    state[i] = THINKING;

    char buf[128];
    snprintf(buf, sizeof(buf), "regresa los cubiertos (%d, %d)",
             i + 1, RIGHT(i) + 1);
    log_event(i, round, buf);

    test(LEFT(i));
    test(RIGHT(i));

    sem_post(&mutex);
}

static void *philosopher(void *arg) {
    int i = *(int *)arg;

    for (int r = 1; r <= ROUNDS; r++) {
        log_event(i, r, "piensa...");
        usleep((useconds_t)THINK_US);

        take_forks(i, r);
        usleep((useconds_t)EAT_US);

        put_forks(i, r);
    }

    printf("[%6ld ms] F[%d] FIN: completo %d rondas.\n",
           now_ms_since_start(), i + 1, ROUNDS);
    fflush(stdout);

    return NULL;
}

int main(void) {
    clock_gettime(CLOCK_MONOTONIC, &t0);

    printf("=== Filosofos comensales | N=%d | RONDAS=%d | THINK_US=%d | EAT_US=%d ===\n\n",
           N_PHILOS, ROUNDS, THINK_US, EAT_US);
    fflush(stdout);

    if (sem_init(&mutex, 0, 1) != 0) {
        perror("sem_init(mutex)");
        return 1;
    }

    for (int i = 0; i < N_PHILOS; i++) {
        state[i] = THINKING;
        if (sem_init(&s[i], 0, 0) != 0) {
            perror("sem_init(s[i])");
            return 1;
        }
    }

    for (int i = 0; i < N_PHILOS; i++) {
        ids[i] = i;
        if (pthread_create(&th[i], NULL, philosopher, &ids[i]) != 0) {
            perror("pthread_create");
            return 1;
        }
    }

    for (int i = 0; i < N_PHILOS; i++) pthread_join(th[i], NULL);

    for (int i = 0; i < N_PHILOS; i++) sem_destroy(&s[i]);
    sem_destroy(&mutex);

    printf("\n=== Listo: todos los filosofos completaron %d rondas. ===\n", ROUNDS);
    return 0;
}
\end{lstlisting}

\newpage
\subsection{Evidencia de Ejecución}
\label{subsec:evidencia}

Para la evidencia se configuró \texttt{ROUNDS = 1} con el fin de que la salida se visualice completa y legible en una sola captura.

\begin{figure}[!htbp]
  \centering
  \fitimage{\detokenize{salidaTerminal.png}}
  \caption{Salida de la terminal (una ronda).}
\end{figure}

\newpage
\subsection{Análisis de Resultados}
\label{subsec:analisis}

Se observa que:

\begin{itemize}
  \item Todos los filósofos completan la ejecución correctamente (en la evidencia: \texttt{1} ronda).
  \item No se presentan bloqueos permanentes (\textit{deadlock}).
  \item La ejecución muestra intercalado de eventos, evidenciando concurrencia real.
  \item Los \textit{timestamps} permiten verificar que no vecinos pueden comer en intervalos cercanos.
\end{itemize}

\newpage
\subsection{Diagramas}
\label{subsec:diagramas}

\subsubsection{Diagrama de Estados del Filósofo}
\label{subsubsec:diag-estados}

\begin{figure}[!htbp]
  \centering
  \fitimage{\detokenize{Diagrama de Estados del Filósofo.png}}
  \caption{Diagrama de estados del filósofo.}
\end{figure}

\newpage
\subsubsection{Diagrama de Recursos y Semáforos}
\label{subsubsec:diag-recursos}

\begin{figure}[!htbp]
  \centering
  \fitimage{\detokenize{Diagrama de Recursos y Semáforos.png}}
  \caption{Recursos compartidos y estructuras de control (mutex y semáforos privados).}
\end{figure}

\newpage
\subsubsection{Diagrama de Sincronización (Lógica del Algoritmo)}
\label{subsubsec:diag-sync}

\begin{figure}[!htbp]
  \centering
  \fitimage{\detokenize{Diagrama de Sincronización - Lógica del Algoritmo.png}}
  \caption{Flujo lógico de sincronización con \texttt{wait}, cambio de estado y \texttt{test}.}
\end{figure}

\newpage
\subsection{Conclusión}
\label{subsec:conclusion}

Se implementó correctamente el problema de los Filósofos Comensales utilizando semáforos POSIX.
La solución evita interbloqueo, permite ejecución concurrente y sincroniza correctamente los accesos a recursos.
Además, es configurable mediante macros para cambiar el número de filósofos y el número de rondas.

\end{document}